\section{Rigorous Convergence Analysis of H-SFP}
\label{app:appendix_convergence}

The convergence dynamics of H-SFP fundamentally diverge from classical federated paradigms. Because the computational tiers are decoupled, optimization at the upper tiers is inherently tied to a sequence of shifting probability measures. Furthermore, we must account for the divergence introduced by asynchronous sub-model aggregation intervals across the hierarchy. We establish convergence through a bottom-up hierarchical induction over the filtration $\mathcal{F}_t$.

Let $(\Omega, \mathcal{F}, \mathbb{P})$ be the underlying probability space, and let $\{\mathcal{F}_t\}_{t \ge 0}$ denote the natural filtration generated by the sequence of stochastic mini-batch realizations. We impose the following regular conditions.

\noindent\textit{Assumption 3 (Fréchet Differentiability and L-Smoothness).} The loss function is $L$-smooth with the Lipschitz constant $0 < L < \infty$. For any $W_1, W_2$ we have
\begin{equation}
\mathcal{L}(W_1) \le \mathcal{L}(W_2) + \langle \nabla \mathcal{L}(W_2), W_1 - W_2 \rangle + \frac{L}{2} ||W_1 - W_2||_2^2.
\end{equation}

\noindent\textit{Assumption 4 (Bounded Variance and Second Moments).} The variances of stochastic gradients for client-side and server-side models across each layer are bounded by $\sigma_l^2$. The expectation of the squared $l_2$-norm of the gradients for these models across each layer is bounded by second-order moments $G_l^2$.

\noindent\textit{Assumption 5 ($L_f$ Lipschitz Feature Manifold).} The feature mappings $W_c$ and $W_e$ are $L_f$ Lipschitz continuous with respect to their parameter space $\Theta$ such that $||W(x; \theta_1) - W(x; \theta_2)||_2 \le L_f ||\theta_1 - \theta_2||_2$.

\noindent\textit{Lemma 1 (Sub-Model Aggregation Divergence).} Before analyzing the distributional drift, we must bound the parameter divergence caused by periodic aggregation. Under the smoothness assumption, the expected squared distance between the aggregated model $\overline{W}_m^t$ and the local model $W_{m,n}^t$ is bounded by the aggregation interval $I_m$
\begin{equation}
\mathbb{E}[||\overline{W}_m^t - W_{m,n}^t||_2^2] \le \mathbb{I}_{\{I_m>1\}} 4\eta^2 I_m^2 \sum_{l=L_{m-1}+1}^{L_m} G_l^2.
\end{equation}
This follows directly from unrolling the weight updates over $I_m$ steps and applying the bounded second moments assumption. The indicator function ensures this bound only applies when tiers delay synchronization ($I_m > 1$).

\noindent\textit{Lemma 2 (Bounding the Edge Tier Distributional Drift).} The edge model $W_e$ optimizes over $\mathbb{Q}_e^t$, a synthesized measure parameterized by the client statistics. Given the measure-theoretic Lipschitz constraint, the training trajectory of $W_e$ absorbs both the weight divergence from Lemma 1 and the temporal distributional drift, converging to a stationary neighborhood. We evaluate the conditional expected progress of the functional $\mathcal{L}_m$ between rounds $t$ and $t+1$. This shift decomposes into the intrinsic optimization descent and the extrinsic distributional drift $\delta_m^t$
\begin{equation}
\mathbb{E}[\mathcal{L}_m^{t+1}(W_e^{t+1})] = \mathbb{E}[\mathcal{L}_m^t(W_e^{t+1})] + \mathbb{E}[\mathcal{L}_m^{t+1}(W_e^{t+1}) - \mathcal{L}_m^t(W_e^{t+1})].
\end{equation}
By our established constraint, the drift is bounded by the Wasserstein distance $\delta_m^t \le L_{stat} \cdot \mathcal{W}_2(\mathbb{Q}_e^{t+1}, \mathbb{Q}_e^t)$. Because the client parameters $W_c$ strictly bound their divergence and stabilize, the generated statistics defining $\mathbb{Q}_e$ also stabilize asymptotically, guaranteeing $\lim_{t \to \infty} \delta_m^t = 0$. Integrating $\delta_m^t$ into the $L$-smoothness expansion and taking the global sum bounds the gradient norm entirely by the intrinsic stochastic noise variance $\sigma_e^2$ and the aggregation drift $\mathbb{I}_{\{I_m>1\}} 4L^2\eta^2 I_m^2 \sum G_l^2$.

\noindent\textit{Proof of Theorem 1 (Global Convergence).} The global classifier operates over the cloud tier synthesized measure $\mathbb{Q}_g^t$. Because the edge tier parameters converge to a stationary neighborhood, their corresponding output feature statistics rigorously stabilize. Consequently, the cloud-level measure generated from these edge prototypes asymptotically approaches a stationary distribution, driving the final global distributional drift $\delta_g^t \to 0$. By substituting the bounded parameter divergence into the smoothness inequality, we can absorb the diminishing distributional drift. Summing over $R$ rounds and dividing by $R$ explicitly yields the bound in Theorem 1, confirming the end-to-end topological convergence of the decoupled H-SFP architecture.